\section*{Bab \ref{ch:g7:how-animals-breathe} --- Bagaimana Hewan Bernapas}

\begin{solution}{exo:g7:how-animals-breathe:1}
Sebuah permukaan pertukaran yang luas, tipis dan lembap, tempat \omterm{prop:g4:journey-of-food:absorb}{darah}
(atau cairan tubuh) bertemu sekelilingnya. Luas: agar penyeberangan
seluruhnya cukup; tipis: agar penyeberangannya mudah; lembap: karena
gasnya menyeberang dalam keadaan terlarut.
\end{solution}

\begin{solution}{exo:g7:how-animals-breathe:2}
\omterm{def:g7:how-animals-breathe:four}{Paru-paru}: \omterm{prop:g3:sorting-living-things:groups}{mamalia}, \omterm{prop:g3:sorting-living-things:groups}{burung} (juga \omterm{prop:g4:classifying-animals:classes}{reptil}, \omterm{prop:g3:sorting-living-things:groups}{amfibi} dewasa). \omterm{def:g7:how-animals-breathe:four}{Insang}: \omterm{prop:g3:sorting-living-things:groups}{ikan},
\omterm{ex:g2:animals-grow-up:frog}{berudu} (juga kerang, larva air). \omterm{def:g7:how-animals-breathe:four}{Trakea}: belalang, lebah --- \omterm{prop:g3:sorting-living-things:groups}{serangga}
pada umumnya. \omterm{def:g7:how-animals-breathe:four}{Pernapasan kulit}: cacing \omterm{def:g6:decomposers-and-soil:soil}{tanah}; \omterm{prop:g3:sorting-living-things:groups}{amfibi} sebagai tambahan.
\end{solution}

\begin{solution}{exo:g7:how-animals-breathe:3}
Air masuk di mulutnya, didorong melewati \omterm{def:g7:how-animals-breathe:four}{insang} yang berumbai ---
tempat pertukarannya terjadi --- lalu keluar di bawah tutup \omterm{def:g7:how-animals-breathe:four}{insangnya}:
sebuah arus satu arah, yang dipompa mulut dan tutupnya secara
bergantian.
\end{solution}

\begin{solution}{exo:g7:how-animals-breathe:4}
Lewat \omterm{def:g7:how-animals-breathe:four}{trakea}: sebuah jaringan pipa udara halus yang berjalan dari
lubang di sepanjang tubuhnya lalu bercabang ke setiap organ ---
udaranya diantar dari pintu ke pintu, sehingga \omterm{prop:g4:journey-of-food:absorb}{darahnya} hanya perlu
membawa sedikit.
\end{solution}

\begin{solution}{exo:g7:how-animals-breathe:5}
Di luar air, permukaan \omterm{def:g7:how-animals-breathe:four}{insangnya} yang berumbai ambruk lalu saling
melekat: permukaannya yang luas menyusut menjadi nyaris tidak ada, dan
tanpa permukaan tidak ada penyeberangan --- \omterm{def:g7:what-is-respiration:respiration}{oksigennya} berlimpah,
mesinnya yang tidak ada untuk mengambilnya.
\end{solution}

\begin{solution}{exo:g7:how-animals-breathe:6}
\omterm{def:g7:how-animals-breathe:four}{Pernapasan kulit} --- dan \omterm{def:g6:exploring-our-environment:conditions}{kondisi} lembapnya yang gagal: \omterm{def:g1:the-five-senses:sense}{kulit} yang
mengering tidak lagi dapat melewatkan gas, dan cacingnya tercekik di
udara terbuka. Luas dan tipis tidak menyelamatkan sebuah permukaan
yang tidak lagi lembap.
\end{solution}

\begin{solution}{exo:g7:how-animals-breathe:7}
\omterm{ex:g2:animals-grow-up:frog}{Berudu}: \omterm{def:g7:how-animals-breathe:four}{insang}. Dewasa: \omterm{def:g7:how-animals-breathe:four}{paru-paru} yang sederhana ditambah \omterm{def:g1:the-five-senses:sense}{kulit} yang
lembap --- dengan \omterm{def:g1:the-five-senses:sense}{kulitnya} sendirian melayani sepanjang musim dingin
yang lesu di dasar kolam.
\end{solution}

\begin{solution}{exo:g7:how-animals-breathe:8}
Mesinnya adalah \omterm{def:g7:how-animals-breathe:four}{paru-paru} --- yang diwarisi dari nenek moyang darat
yang memberi \omterm{def:g2:animals-grow-up:born}{susu} --- sehingga ia harus muncul ke permukaan untuk
bernapas dengan udara, semirip apa pun kehidupannya dengan \omterm{prop:g3:sorting-living-things:groups}{ikan}: kalau
ditahan di bawah, ia tenggelam. Pakaiannya milik airnya; mesinnya
milik keluarganya.
\end{solution}

\begin{solution}{exo:g7:how-animals-breathe:9}
Sebuah simpanan udara yang dibawa --- tangki seorang penyelam di bawah
tutup sayapnya, yang darinya \omterm{def:g7:how-animals-breathe:four}{trakeanya} menimba. Ia harus kembali ke
permukaan secara teratur untuk memperbarui gelembungnya.
\end{solution}

\begin{solution}{exo:g7:how-animals-breathe:10}
(a) Kerang: penghuni air, tanpa gerakan yang terlihat tetapi
berinsang di dalam \omterm{def:g1:animals-around-us:covering}{cangkangnya}, yang dicuci arus yang dipompanya. (b)
Belalang: penghuni udara; deretan lubang mungil, perut yang berdenyut
--- \omterm{def:g7:how-animals-breathe:four}{trakea}. (c) Ular rumput: seorang \omterm{prop:g4:classifying-animals:classes}{reptil} yang pulang pergi ---
\omterm{def:g7:how-animals-breathe:four}{paru-paru}; ia berenang dengan lubang hidung di atas lalu muncul ke
permukaan untuk bernapas.
\end{solution}

\begin{solution}{exo:g7:how-animals-breathe:11}
\omterm{def:g7:how-animals-breathe:four}{Trakea} mengantar lewat pipa udara yang bercabang, hebat pada jarak
milimeter dan terlalu lambat pada jarak yang panjang: bagian dalam
sebuah tubuh sebesar paus akan terletak terlalu jauh dari lubang mana
pun. Mesinnya menetapkan langit-langit ukuran, dan \omterm{prop:g3:sorting-living-things:groups}{serangga} hidup di
bawahnya.
\end{solution}

\begin{solution}{exo:g7:how-animals-breathe:12}
Spesifikasinya: sebuah permukaan pertemuan yang luas, tipis dan lembap
antara tubuh dan sekelilingnya. \omterm{def:g7:how-animals-breathe:four}{Insang} memenuhinya dengan permukaan
berumbai yang direntangkan airnya; \omterm{def:g7:how-animals-breathe:four}{paru-paru} dengan jutaan kantong
yang terlipat dalam naungan yang lembap; \omterm{def:g7:how-animals-breathe:four}{trakea} dengan jumlah
permukaan dalam ujung terhalusnya, yang lembap pada ujungnya; \omterm{def:g1:the-five-senses:sense}{kulit}
dengan seluruh bagian luar tubuhnya yang lembap, yang dijaga tetap
luas dibandingkan tubuh yang kecil dan tipis. Empat pemegang, satu hal
yang dipegang.
\end{solution}

\begin{solution}{pb:g7:how-animals-breathe:1}
\textbf{1.} Benih \omterm{prop:g3:sorting-living-things:groups}{ikan} trout dan \omterm{prop:g3:sorting-living-things:groups}{ikan} duri: \omterm{def:g7:how-animals-breathe:four}{insang}. \omterm{ex:g2:animals-grow-up:frog}{Berudu}: \omterm{def:g7:how-animals-breathe:four}{insang}
(yang sudah berkaki sedang di tengah pembangunan ulang). Katak dewasa:
\omterm{def:g7:how-animals-breathe:four}{paru-paru} ditambah \omterm{def:g1:the-five-senses:sense}{kulit}. Siput kolam: \omterm{def:g7:how-animals-breathe:four}{paru-paru} --- pengunjung
permukaan. Kumbang selam: \omterm{def:g7:how-animals-breathe:four}{trakea}, yang diberi makan dari gelembung
yang dibawanya. Larva lalat sehari: \omterm{def:g7:how-animals-breathe:four}{insang} (berkas yang berumbai).
Cacing \omterm{def:g6:decomposers-and-soil:soil}{tanah}: \omterm{def:g1:the-five-senses:sense}{kulit}.

\textbf{2.} \omterm{def:g7:how-animals-breathe:four}{Insang}. Lambaiannya memperbarui air yang bersentuhan
dengan permukaan \omterm{def:g7:how-animals-breathe:four}{insangnya} --- sebuah arus buatan sendiri di tempat
kolam yang tenang tidak menyediakannya, sehingga air segar pembawa
\omterm{def:g7:what-is-respiration:respiration}{oksigen} tetap ada di penyeberangannya.

\textbf{3.} Siput kolamnya, kumbang selamnya, katak dewasanya (dan
ular rumput mana pun yang berkunjung): semua pemakai \omterm{def:g7:how-animals-breathe:four}{paru-paru} dan
gelembung --- mesinnya mengambil \omterm{def:g7:what-is-respiration:respiration}{oksigen} dari udara, sehingga
permukaannya adalah pom pengisiannya.

\textbf{4.} Pembangunan ulang \omterm{def:g7:what-is-respiration:respiration}{pernapasan} pada \omterm{def:g2:animals-grow-up:metamorphosis}{metamorfosisnya}: \omterm{def:g7:how-animals-breathe:four}{insang}
yang melayani \omterm{ex:g2:animals-grow-up:frog}{berudu} tanpa kaki sedang ditukar dengan \omterm{def:g7:how-animals-breathe:four}{paru-paru} (dan
\omterm{def:g1:the-five-senses:sense}{kulit} yang bernapas) yang akan dibutuhkan katak kecil yang menuju
tepian.

\textbf{5.} Para pemakai \omterm{def:g7:how-animals-breathe:four}{insang} menderita lebih dahulu: seluruh
pasokannya adalah persediaan terlarut airnya yang menyusut. Para
pengunjung permukaan --- siput, kumbang, katak --- menimba udara tanpa
batas di atasnya dan sekadar pulang pergi lebih sering.

\textbf{6.} Milimeter teratasnya bersentuhan dengan udara: \omterm{def:g7:what-is-respiration:respiration}{oksigen}
melarut masuk dari atas, sehingga lapisan selaputnya adalah air yang
paling tidak miskin di sebuah kolam yang tercekik. \omterm{prop:g3:sorting-living-things:groups}{Ikan} durinya sedang
mengantre di sumur yang terakhir.

\textbf{7.} \omterm{def:g7:how-animals-breathe:four}{Insangnya} tersumbat: lumpur halus mengendap pada permukaan
berumbainya lalu menghalangi penyeberangannya. Harapkanlah para
pernapas \omterm{def:g7:how-animals-breathe:four}{insang} menghindari awannya, dan pemompaan tutupnya
kepayahan --- dengan kejang batuk untuk membersihkan endapannya.

\textbf{8.} \omterm{def:g7:how-animals-breathe:four}{Pernapasan kulit}, dan \omterm{def:g6:exploring-our-environment:conditions}{kondisi} lembap-tetapi-tanpa-udaranya:
liangnya yang tergenang memuat terlalu sedikit \omterm{def:g7:what-is-respiration:respiration}{oksigen}, sehingga
cacingnya muncul ke permukaan untuk bernapas --- menukar tercekik di
bawah dengan risiko mengering di atas.

\textbf{9.} Tabelnya seperti yang ditugaskan (mesin dan medium tiap
\omterm{def:g1:animals-around-us:animal}{hewannya}; yang pulang pergi: siput, kumbang, katak). Yang mengepalai
setiap lajurnya: permukaan pertukaran yang luas-tipis-lembap ---
satu spesifikasi yang dipegang keempat mesinnya.

\textbf{10.} Katak: tercukupi --- lesu di dalam lumpurnya, dengan
\omterm{def:g1:the-five-senses:sense}{kulitnya} yang bernapas melayani kebutuhan musim dinginnya yang kecil.
Siput dan kumbang: kesulitan --- permukaannya tersegel; mereka
melewati musim dingin dalam keadaan dorman dengan kebutuhan yang
diturunkan atau dengan gelembung yang terperangkap. \omterm{prop:g3:sorting-living-things:groups}{Ikan}: tercukupi
selama persediaan airnya bertahan --- dingin memperlambat setiap bunyi
mesin (\cref{ex:g7:what-is-respiration:demand}), dan \omterm{def:g7:how-animals-breathe:four}{insangnya}
melayani permintaan yang berkurang di bawah esnya.

\textbf{11.} Lonceng sutra laba-labanya adalah tangki udara yang
tetap, sedangkan gelembung kumbangnya adalah tangki yang dibawa:
keduanya menyimpan udara permukaan di bawah air dan keduanya tetap
melayani sebuah mesin udara (perlengkapan \omterm{def:g7:how-animals-breathe:four}{paru-paru} buku dan \omterm{def:g7:how-animals-breathe:four}{trakea}
milik laba-labanya) di dalam medium yang keliru.

\textbf{12.} Karena kolamnya menawarkan \omterm{def:g7:what-is-respiration:respiration}{oksigen} dalam dua medium ---
yang langka dan terlarut di bawah, yang kaya di atas --- dan tiap
mesinnya memanen satu medium pada satu ukuran kehidupan: hanya dengan
keempatnya bekerja, setiap sudutnya, dari lumpur sampai selaput
permukaannya, menjadi terhuni.
\end{solution}
